% 文件名保留兼容；按 role=evaluation 定位，三问为第7章，其他问数自动编号。
% 正常字号/行距/页边距下净正文至少一整页，允许自然跨页。
% 以下仅为待填写骨架；空条目不能通过生成与审计，必须替换为本题真实评价。
\section{模型评价与展望}

\subsection{模型优点}

\begin{enumerate}
  \item % 写明确优点及具体解释，对应前文真实方法或验证依据。
  \item % 写另一项独立优点及其依据，不重复第一点。
  \item % 至少第三项独立优点及其依据，可按实际需要增加。
\end{enumerate}

\subsection{模型缺点}

\begin{enumerate}
  \item % 说明具体局限、原因及对结果或适用范围的影响。
  \item % 说明另一项独立局限、原因及其影响。
  \item % 至少第三项独立局限、原因及其影响，可按实际需要增加。
\end{enumerate}

\subsection{模型的推广与改进}

% 第一自然段：阐述可迁移场景、适用条件与边界，不写成分点列表。

% 第二自然段：针对上述局限提出具体改进路径，区分已验证结论和未来设想。
% 至少两段，以空行或 \par 分段，不用强制换行、手写编号替代自然段。
% 最终实际查看本章全部PDF页面及测量明细，不能只凭页码、字数或工具PASS。
